\subsection{Full simulation: origin of the low-momentum excess}
\label{sec:sim}

The SE excess over the ME between the p and d peaks at low momentum is a fixed rate per track, flat in centrality and flat in time
up to 27\,ns (Sec.~\ref{sec:lowp}). Two track-correlated sources were suspected: decays in flight and particles scattered or produced
in the detector material. Both are present in a full simulation, where every BTOF hit can be traced to the particle that made it.

\paragraph{Simulation set-up.}
\begin{itemize}
\item Generator: HIJING 1.383, Au+Au at $\sqrt{s_{NN}}=19.6$\,GeV, impact parameter 0--14.5\,fm, no jet quenching, weak decays
  ($\pi^0$, K$^0_S$, $\Lambda$, \ldots) left to GEANT (\fpath{sim/starsim_hijing.C}). All particles go to GEANT (no $\eta$ cut).
  Vertex Gaussian with $\sigma_{x,y}=0.1$\,cm, $\sigma_z=50$\,cm.
\item Detector: starsim (GEANT3) with the y2019a AgML geometry (iTPC, BTOF, VPD, magnet).
\item Reconstruction: \fpath{sim/reco_sim.C} with the bfc options \fcmd{y2019a}, \fcmd{fzin}, \fcmd{AgML}, \fcmd{useXgeom},
  \fcmd{iTpcIT}, \fcmd{MakeEvent}, \fcmd{ITTF}, \fcmd{VFMinuit}, \fcmd{tpcDB}, \fcmd{btofsim}, \fcmd{vpdSim}, \fcmd{btofMatch},
  \fcmd{btofCalib}, \fcmd{vpdCalib}, \fcmd{IdTruth}, \fcmd{McEvent}, \fcmd{CMuDst} (full list in \fpath{sim/sim_job.sh}), plus either
  \fcmd{tfs} (fast TPC simulator) or \fcmd{TpcRS}, \fcmd{TpxClu}, \fcmd{TpcHitMover} (full TPC response). The BTOF matching and calibration are the production makers.
\item BTOF simulator: the SL24y \fpath{StBTofSimMaker} in fast mode, with a local patch (\fpath{StRoot/StBTofSimMaker}). The SL24y
  maker sets the truth id of a hit only for crosstalk hits; the patch sets it for every hit to the GEANT track of the earliest
  contribution in the cell, i.e.\ the one that defines the hit time.
\item Truth analysis: \fpath{StRoot/StTofSimAnaMaker} (\fpath{macros/runTofSimAna.C}) on the simulation MuDst. Same event and track
  selection as the data (Table~\ref{tab:trkcut}). The matched hit of every TOF-matched QA track is classified with the MC record
  (\fcmd{StMuMcTrack}, \fcmd{StMuMcVertex}):
  \emph{own} (the track's particle made the hit); \emph{decay} (a descendant through a decay: weakly decaying parent and a muon daughter,
  or a decay vertex at $R<190$\,cm); \emph{secondary} (a descendant through any other process, mostly interactions); \emph{relative}
  (an ancestor or sibling); \emph{unrelated} (no common ancestry, the part a mixed event describes).
  $m^2$ is computed with the true event time (hit time minus the simulation clock offset of 19.01\,ns).
\item Samples: \simEvGen{} events with the fast TPC simulator (999 jobs of 50 events; one job hung in HIJING and was removed),
  \simEvAna{} after the vertex cuts, \simNTrk{} QA tracks; and the same \simEvGenOneK{} GEANT events reconstructed both with the fast
  simulator and with TpcRS. Production scripts \fpath{sim/sim_job.sh}, \fpath{sim/sim_job2.sh}, \fpath{sim/sim.sub},
  \fpath{sim/sim2.sub}; numbers from \fpath{scripts/make_sim.C} (\fpath{logs/sim_numbers.txt}).
\end{itemize}

\choice{The fast TPC simulator is used for the large sample. The full TPC response costs about 3 minutes per peripheral and 10 minutes
per mid-central event (98\% of the CPU), against a few seconds. On the same 1k events the two give the same fractions of all match
categories within \simTrsMaxSig\,$\sigma$ (Table~\ref{tab:simtrs}).}

\input{tables/sim_trs}

\caveat{Limitations of the simulation: (i) the fast BTOF simulator gives every hit ToT$\,=0$, so the production rule ``largest ToT
wins'' between several candidate hits cannot act and ties go to the next criterion; in data a slow, heavily ionising secondary has a
large ToT and is preferred. (ii) Production of slow protons and nuclear fragments in the material depends on the GEANT3 hadronic
package and its cuts (not tuned here). (iii) Centrality: HIJING plus reconstruction was not tuned to the data refMult; simulated
events are classified by their refMult percentile within the simulated sample (b\,$<$\,14.5\,fm), an approximate match to the
data classes. (iv) HIJING has no light nuclei (no deuteron peak). (v) No noise hits and no pile-up.}

\subsubsection{Where the matched hits come from}
\paragraph{Track selection and the decay category.} The simulated tracks pass exactly the data selection: primary tracks with
$p_T>0.15$\,GeV/$c$, $|\eta|<1$, at least 15 fit points, fit/possible points $>0.52$, global DCA $<3$\,cm, and the production BTOF matching
(\fpath{StTofSimAnaMaker}, same cuts as \fpath{StTofMEOverlayMaker}). There is no explicit requirement on where a decay happens. The MC
particle of a TPC track is the one that contributed most of its hits (IdTruth). This fixes which decays can appear in the decay category:
\begin{itemize}
\item A $\pi/K\to\mu\nu$ decay early in the TPC gives a track whose hits are mostly from the muon; its IdTruth is the muon, so the TOF hit
  of that muon is counted as \emph{own}. At 0.2--0.4\,GeV/$c$, \ownMuA\% of the own hits are made by muons (decays before the outer TPC
  and the few primary muons). Decays very close to the vertex mostly fail the DCA and fit-point cuts.
\item The decay category therefore contains decays late in the TPC or after it: only \decRALo\% of the decay vertices are at $R<100$\,cm;
  \decRATpc\% are at 100--190\,cm (outer TPC) and \decRAOut\% between the TPC and the TOF, in the TOF trays included, at 0.2--0.4\,GeV/$c$
  (\decRBTpc\%/\decRBOut\% at 0.4--0.6 and \decRCTpc\%/\decRCOut\% at 0.6--1\,GeV/$c$). The track is reconstructed with the parent's
  momentum and the hit comes from the muon, which keeps nearly the same direction and velocity.
\item These hits stay at the pion mass: \decPiA\%, \decPiB\% and \decPiC\% of them have $|q\,m^2|<0.15$\,GeV$^2$/$c^4$ in the three
  momentum bins (Fig.~\ref{fig:sim_qm2}); the peak is slightly broader than the own-pion peak, with a tail toward the kaon at 0.6--1\,GeV/$c$.
\end{itemize}
The same decays happen in the data: they are part of the pion peak of the SE and are not in the ME. They are true matches of a
pion (or kaon) track with a hit of its own decay product at almost the expected time, so they count as signal, not background, and they do
not contribute to the late p--d excess.

Table~\ref{tab:simfrac} gives the categories per momentum bin. At 0.2--0.4\,GeV/$c$, \simFrOwnA\% of the matches are own hits,
\simFrDecA\% decay muons, \simFrSecA\% secondaries and \simFrUnrA\% unrelated hits.
Fig.~\ref{fig:sim_qm2} shows how they distribute in $q\times m^2$ (read by eye): decay muons stay close to the pion peak; secondaries
that arrive together with the track's own particle sit at its $m^2$ (pion and proton peaks), and the late ones form a broad
continuum; at 0.2--0.4\,GeV/$c$ this continuum is above the unrelated matches for $|q\,m^2|>1$.

\input{tables/sim_frac}

\begin{figure}[H]
\centering
\figpanel{0.9}{sim_qm2}{$q\times m^2$ of TOF-matched QA tracks in HIJING+GEANT (\simEvAna{} events, all centralities, fast TPC simulator),
by origin of the matched hit: own particle (gray), decay (green), secondary (red), unrelated (blue). The relative category is empty.
0.04\,GeV$^2$/$c^4$ bins, log scale.}
\caption{Simulation: origin of the matched hit versus $q\times m^2$ (\fpath{scripts/make_sim.C}).}
\end{figure}

Fig.~\ref{fig:sim_origin} shows where the branch leading to the hit starts. \simSecTray\% of the secondaries (0.2--1\,GeV/$c$) start at
$200<R<240$\,cm, i.e.\ in the BTOF trays themselves; the decays are spread along the path from $R\approx110$\,cm to the trays (read by eye). The
secondaries that make the hit are e$^\pm$ (\simSecElec\%), $\pi^\pm$ (\simSecPion\%) and protons (\simSecProton\%); the decay hits are muons.

\begin{figure}[H]
\centering
\figpanel{0.95}{sim_origin}{Left: radius of the vertex where the branch leading to the hit leaves the track's particle, for decays (green)
and secondaries (red), $0.2<p<1$\,GeV/$c$, log scale. Right: particle that made the hit, in \% of each category.}
\caption{Simulation: origin of correlated hits (\fpath{scripts/make_sim.C}).}
\end{figure}

\subsubsection{Comparison with the data}
In the p--d window ($1.6<|q\,m^2|<2.8$\,GeV$^2$/$c^4$) at 0.2--0.4\,GeV/$c$ the simulation has \simHiNCorr{} correlated entries;
\simHiSecShare\% of them are secondaries, \simHiNDec{} decays, and only \simHiNOwn{} entries are own hits.
Fig.~\ref{fig:sim_rate} and Table~\ref{tab:simrate} compare, per $10^4$ QA tracks and versus centrality, the data SE$-$ME (method 2)
with the simulated correlated matches, and the data ME with the simulated unrelated matches.
\begin{itemize}
\item 0.2--0.4\,GeV/$c$: the simulated correlated rate is \simCorrA{} per $10^4$ tracks (all centralities) and, like the data, flat in
  centrality. The data/simulation ratio is \simDataRatMinA--\simDataRatMaxA{} across the five classes, with at most
  \simNsigMaxA\,$\sigma$ difference.
\item 0.4--0.6\,GeV/$c$: the simulation (\simCorrB) is below the data in every class; data/simulation \simDataRatMinB--\simDataRatMaxB,
  up to \simNsigMaxB\,$\sigma$.
\item 0.6--1\,GeV/$c$: the simulation (\simCorrC) agrees with the data, at most \simNsigMaxC\,$\sigma$ apart.
\item The simulated unrelated rate rises with centrality like the data ME but stays below it (Table~\ref{tab:simrate}): the simulation
  has fewer BTOF hits per event (no noise, no pile-up, untuned multiplicity).
\end{itemize}

\begin{figure}[H]
\centering
\figpanel{0.98}{sim_rate}{Entries in $1.6<m^2<2.8$\,GeV$^2$/$c^4$ per $10^4$ QA tracks versus centrality, for 0.2--0.4, 0.4--0.6 and
0.6--1\,GeV/$c$. Black full: data SE$-$ME (method 2); red squares: simulated correlated matches (decay + secondary + relative); black
open, dashed: data ME; blue open, dashed: simulated unrelated matches. Simulation centrality from the refMult percentile (approximate).
Statistical errors.}
\caption{Simulation versus data: the track-correlated rate in the p--d window (\fpath{scripts/make_sim.C}).}
\end{figure}

\input{tables/sim_rate}

\paragraph{Time structure.} Fig.~\ref{fig:sim_dt} compares $t-t_\pi$ per QA track. In the late region 20--27\,ns, beyond the K and p
peaks, the data SE$-$ME is \lateDataA, \lateDataB{} and \lateDataC{} per ns per $10^4$ tracks at 0.2--0.4, 0.4--0.6 and 0.6--1\,GeV/$c$;
the simulated correlated matches give \lateSimA, \lateSimB{} and \lateSimC{} (data/simulation \lateRatA, \lateRatB, \lateRatC). At the
lowest momentum the simulation has the late plateau at a similar level, but not its momentum independence: in data the plateau stays at the
same level up to 1\,GeV/$c$, in the simulation it falls with momentum.

\begin{figure}[H]
\centering
\figpanel{0.98}{sim_dt}{$t-t_\pi$ per ns per $10^4$ QA tracks. Black points: data SE$-$ME (method 2, 0--80\%); gray: simulation, own +
correlated matches (everything except unrelated); red: simulated correlated matches. Simulation with the true event time; the peak
widths and the K/p yields of HIJING are not tuned to the data. Log scale.}
\caption{Simulation versus data: time structure (\fpath{scripts/make_sim.C}).}
\end{figure}

\paragraph{Charge.} Table~\ref{tab:simcharge} compares the negative/positive ratio. The data excess is positive-dominated at
0.2--0.3\,GeV/$c$ (\chDataA) and below 1 elsewhere; the simulated correlated matches are charge-symmetric at 0.2--0.3\,GeV/$c$
(\chSimA) and negative-dominated at 0.3--0.6\,GeV/$c$ (\chSimB, \chSimC). The positive-only data component at the lowest momentum
(Sec.~\ref{sec:lowp}) is therefore not reproduced (HIJING has no light nuclei, and the proton energy-loss tail is small in the
simulation), and the negative excess of the simulation at 0.3--0.6\,GeV/$c$ is not seen in data.

\input{tables/sim_charge}

\subsubsection{Data mass spectrum with the simulated correlated component}
\label{sec:simtmpl}
The simulated correlated matches (decay + secondary + relative), normalized per QA track to the data, are added to the data ME and
compared with the data SE $m^2$ spectrum (Fig.~\ref{fig:simtmpl}, \fpath{scripts/make_simtmpl.C}). The simulated $m^2$ is
$|q\,m^2|$ and its centrality is the refMult percentile of the simulated events.
\begin{itemize}
\item Between the p and d peaks (1.6--2.8\,GeV$^2$/$c^4$) at 0.2--0.4\,GeV/$c$, SE/ME is \tmplFastRatMECiA, \tmplFastRatMEMidA{} and
  \tmplFastRatMEPerA{} in 0--10, 20--40 and 60--80\%; SE/(ME + sim.\ correlated) is \tmplFastRatSumCiA, \tmplFastRatSumMidA{} and
  \tmplFastRatSumPerA. At 0.6--1\,GeV/$c$: \tmplFastRatSumCiC, \tmplFastRatSumMidC, \tmplFastRatSumPerC. At 0.4--0.6\,GeV/$c$ an excess
  remains: \tmplFastRatSumCiB, \tmplFastRatSumMidB, \tmplFastRatSumPerB.
\item In the K--p valley (0.35--0.6\,GeV$^2$/$c^4$) and above the deuteron (3.0--4.5\,GeV$^2$/$c^4$) at 0.2--0.4\,GeV/$c$ the simulation
  predicts far more correlated matches than the data allow: SE/(ME + sim.) is \tmplFastRatSumLoCiA--\tmplFastRatSumLoPerA{} and
  \tmplFastRatSumHiCiA--\tmplFastRatSumHiPerA{} (0--10\% and 60--80\%), i.e.\ the simulated correlated component alone is larger than the
  data SE there (Table~\ref{tab:simvartmpl}).
\end{itemize}

\input{tables/simvar_tmpl}

\begin{figure}[H]
\centering
\figpanel{0.98}{simtmpl_h50k}{Data SE $m^2$ (points) compared with the data ME (method 2, red filled), the simulated correlated matches
(HIJING+GEANT, 50k events, no ToT; green dashed) scaled to the number of data QA tracks, and their sum (blue). Columns: 0--10, 20--40,
60--80\%; rows: 0.2--0.4, 0.4--0.6, 0.6--1\,GeV/$c$. 0.1\,GeV$^2$/$c^4$ bins, log scale. The ratios in the p--d window are printed in
each panel.}
\caption{Data mass spectra with the simulated correlated component (\fpath{scripts/make_simtmpl.C}).}
\label{fig:simtmpl}
\end{figure}

\subsubsection{BTOF response: ToT and the slow simulator}
\label{sec:simvar}
The SL24y BTOF simulator writes the trailing edge equal to the leading edge in both its modes, so every simulated hit has ToT\,$=0$ and
the production rule ``largest ToT wins'' between several candidate hits never acts (all candidates tie and the closest-$y$ rule
decides); the slow mode has no setter. The patched \fpath{StBTofSimMaker} adds, through environment variables: (a) a ToT from the cell
energy deposit in fast mode, ToT$\,=16.0+1.05\ln(dE/dE_{\rm MIP})$\,ns with $dE_{\rm MIP}=0.047$\,keV (the median GEANT deposit of
minimum-ionizing pions), and (b) the slow (avalanche) mode with its threshold, with the ToT from the summed avalanche charge in the same
form ($Q_{\rm MIP}=0.0095$\,pC). The constants are taken from the data: in Run~19 the ToT of identified pions is 16.0\,ns (16--84\%:
15.1--17.1\,ns) at minimum ionization and rises to 17.9\,ns for protons at $\beta\gamma\approx0.45$ (\fpath{macros/totData.C}, one MuDst
file). The same 5k GEANT events (\fpath{sim/sim_job3.sh}) were reconstructed with the three responses (Table~\ref{tab:simvar}).

\input{tables/simvar_tab}

\begin{itemize}
\item With a ToT the matching efficiency of QA tracks rises by 2--3 points (\varNoTotEffA\% $\to$ \varTotEffA\%, \varSlowEffA\% at
  0.2--0.4\,GeV/$c$), and the unrelated matches roughly double to triple: when several hits compete, the ToT rule picks a random hit more
  often than the closest-$y$ rule. At 0.4--0.6\,GeV/$c$ the unrelated rate in the p--d window becomes \varTotPdBU{} (fast + ToT) and
  \varSlowPdBU{} (slow) per $10^4$ tracks, against \varNoTotPdBU{} without ToT; the data ME is \varDataPdBU.
\item The correlated rate in the p--d window at 0.4--0.6\,GeV/$c$ grows from \varNoTotPdB{} to \varTotPdB{} and \varSlowPdB{} per $10^4$
  tracks (data SE$-$ME \varDataPdB), closing part of the gap seen in Table~\ref{tab:simrate}. At 0.2--0.4\,GeV/$c$ it is unchanged within
  errors (\varNoTotPdA, \varTotPdA, \varSlowPdA; data \varDataPdA).
\item The over-prediction in the K--p valley and above the deuteron at 0.2--0.4\,GeV/$c$ is not removed: \varNoTotKpA, \varTotKpA, \varSlowKpA{}
  correlated per $10^4$ tracks against a data SE$-$ME of \varDataKpA{} in the valley, and \varNoTotHiA, \varTotHiA, \varSlowHiA{} against
  \varDataHiA{} at 3.0--4.5\,GeV$^2$/$c^4$. The slow simulator's threshold does not remove these secondaries either.
\end{itemize}

\caveat{The 5k sample has 1--70 entries per cell of Table~\ref{tab:simvar}; differences below about 2$\sigma$ are not significant. The
ToT scale of the simulation is not tuned: own hits at 0.4--1\,GeV/$c$ have a median ToT of about 18\,ns (data about 16\,ns), which matters
only through ToT differences between candidates. In slow mode the avalanche charge is saturated and does not depend on $\beta\gamma$, so
the slow ToT does not favour heavily ionizing hits.}

\choice{Conclusion: the simulation confirms that most of the low-momentum excess is track-correlated and not combinatorial. It is
dominated by secondaries produced by the track's own particle in the BTOF trays (R\,$\approx$\,210--220\,cm), not by decays in flight,
which stay under the pion peak. At 0.2--0.4\,GeV/$c$ the simulated rate is flat in centrality and within \simDataRatMinA--\simDataRatMaxA{}
(data/simulation) of the data.}

\caveat{Not reproduced: the data rate at 0.4--0.6\,GeV/$c$ (data/simulation \simDataRatMinB--\simDataRatMaxB{} without ToT; a ToT
closes part of it, Sec.~\ref{sec:simvar}), the momentum-independent late plateau, the charge pattern, and the size of the correlated
component outside the p--d window: at 0.2--0.4\,GeV/$c$ the simulation predicts about 3--4 times more correlated matches than the data
allow in the K--p valley and above the deuteron, with or without ToT and with the slow simulator. The simulated secondaries in the trays
are therefore not a reliable absolute template; the likely causes are the GEANT3 treatment of low-energy secondaries in the tray
material and the absence of any electronics dead time or noise in the simulation (not tested).}

\choice{Decided (2026-11-06): the simulation is used as evidence for the origin of the low-momentum excess, not as a correction. A
simulated template would fix the p--d window at 0.2--0.4\,GeV/$c$ but over-subtract the K--p valley and the high-$m^2$ region by a factor
of 2--3. The component is estimated from the data instead (Sec.~\ref{sec:corrdd}).}
